\subsection{Neural Networks}

In today's world, where we have an abundance of data and easy access to it, the power of machine learning has become a game changer in many fields. At its core, machine learning relies on algorithms to identify patterns in data and to make decisions with minimal human intervention. Tom Mitchell, a leading figure in the field, explains it as \cite{mitchell_machine_1997}:
\begin{quote}
"A computer program is said to learn from experience E with respect to some task T and some performance measure P, if its performance on T, as measured by P, improves with experience E."
\end{quote}
This captures the essence of machine learning: creating models that automatically improve through experience. Neural networks, the topic of this section, are a class of machine learning models that operate within this framework by simulating the complex networks of the human brain.

\subsubsection{Fundamentals and History}

Neural Networks (NN), named for their similarity to the structure of the human brain, are algorithms designed for pattern recognition. In the human brain, a neuron receives signals through structures called dendrites, processes these signals in the nucleus, and passes the resultant processed signal down its axon to the dendrites of adjacent neurons. Similarly, in a neural network, each node represents an artificial neuron that receives input, processes it through a mathematical function, and passes the output to the subsequent layer of nodes. In both cases, the strength and nature of the connections between neurons determine the behaviour of the network \cite{goodfellow_deep_2016}. Refer to Figure~\ref{Figure:HumanNeuronVsArtificialNeuron} for a visual representation of the concepts explained above.

\begin{figure}[htb!]
\centering
\includegraphics[width=0.7\textwidth]{Images/HumanNeuronVsArtificialNeuron.png}
\caption{Human Neuron compared to an Artificial Neuron \cite{pramoditha_human_nodate}.}
\label{Figure:HumanNeuronVsArtificialNeuron}
\end{figure}

The concept of neural networks dates back to the 1950s with the perceptron, a simple model of a biological neuron developed by Frank Rosenblatt \cite{rosenblatt_perceptron_1958}. The perceptron marked the early stages of NN research, but was limited to linearly separable datasets, leading to some periods of reduced interest in the field, known as "AI winters." However, the 1980s brought a significant breakthrough with the backpropagation algorithm, enabling neural networks to learn patterns in non-linearly separable datasets using complex and multilayered models \cite{rumelhart_learning_1986}.

\subsubsection{Key Concepts and Feed-Forward}

In a neural network, the feed-forward computation within each layer \( l \) is based on the data processed by its predecessor. For the first layer, this input \( x \) is the raw data itself, while for subsequent layers, it is the output of the previous layer \( h^{(l-1)}(x) \). The computation in each layer can be formalised as \cite{goodfellow_deep_2016}:

\begin{equation}
z^{(l)}(x) = W^{(l)} h^{(l-1)}(x) + b^{(l)}
\end{equation}

Here, the weight matrix \( W^{(l)} \in \mathbb{R}^{K_l \times K_{l-1}} \) and the bias vector \( b^{(l)} \in \mathbb{R}^{K_l} \) are required for the computation at layer \( l \), with \( K_l \) representing the number of neurons or units in layer \( l \), and \( K_{l-1} \) indicating the number of neurons in the preceding layer \( l-1 \). In particular, \( W^{(1)} \in \mathbb{R}^{K_1 \times D} \) for the first layer, where \( D \) is the size of the input data. The activation function \( g \) is subsequently applied to the pre-activation \( z^{(l)}(x) \in \mathbb{R}^{K_l} \) to produce the activated output \( h^{(l)}(x) \in \mathbb{R}^{K_l} \) \cite{goodfellow_deep_2016}:

\begin{equation}
h^{(l)}(x) = g(z^{(l)}(x))
\end{equation}

During the design of a neural network, the choice of activation functions is very important. Starting by the hidden layers, we may employ a variety of non-linear activation functions, adding layers of expressiveness to the network and allowing it to learn more complex patterns. Finally, the activation function used in the output layer is determined by the type of task the network is going to perform. This careful selection directly influences the network's ability to accurately and effectively model the dataset \cite{goodfellow_deep_2016}. Refer to Table~\ref{Tables:ActivationFunctions} for a summary of the most common activation functions.

\begin{table}[htb!]
\centering
\footnotesize
\begin{tabularx}{\textwidth}{@{}lXl@{}}
\toprule
\textbf{Activation Function} & \textbf{Formula} & \textbf{Description} \\
\midrule
Linear & \( a(x) = x \) & Suitable for regression tasks. \\
\addlinespace
Sigmoid & \( \sigma(x) = \frac{1}{1 + e^{-x}} \) & Maps inputs to (0, 1), for binary classification. \\
\addlinespace
Hyperbolic Tangent & \( \tanh(x) = \frac{e^{x} - e^{-x}}{e^{x} + e^{-x}} \) & Maps inputs to (-1, 1), offering richer gradients than sigmoid. \\
\addlinespace
ReLU & \( \text{ReLU}(x) = \max(0, x) \) & Encourages sparse activation, helps with gradient flow. \\
\addlinespace
Softmax & \( \text{Softmax}(z)_i = \frac{e^{z_i}}{\sum_{k=1}^K e^{z_k}} \) & Turns logits into probabilities, for multi-class classification. \\
\bottomrule
\end{tabularx}
\caption{Summary of the most common Activation Functions \cite{goodfellow_deep_2016}.}
\label{Tables:ActivationFunctions}
\end{table}

\subsubsection{Learning using Backpropagation}

Neural networks learn by iteratively adjusting weights and biases to minimise the error between predicted outputs and actual target values. This learning process is guided by a loss function that quantifies the model's error. The goal is to optimise the network by finding parameters that minimise this loss function. Different problems require different loss functions, such as mean squared error for regression tasks and cross-entropy for classification \cite{goodfellow_deep_2016}. Refer to Table~\ref{Tables:LossFunctions} for a detailed overview of their formulas.

\begin{table}[htb!]
\centering
\footnotesize
\begin{tabularx}{\textwidth}{@{}lXl@{}}
\toprule
\textbf{Loss Function} & \textbf{Formula} & \textbf{Application} \\
\midrule
Mean Squared Error (MSE) & \( \frac{1}{n} \sum_{i=1}^n (y_i - \hat{y}_i)^2 \) & Regression \\
\addlinespace
Negative Log-Likelihood (Cross-Entropy) & \( -\sum_{i=1}^n y_i \log(\hat{y}_i) \) & Classification \\
\bottomrule
\end{tabularx}
\caption{Summary of the most common Loss Functions \cite{goodfellow_deep_2016}.}
\label{Tables:LossFunctions}
\end{table}

The core algorithm used to train neural networks is known as backpropagation \cite{rumelhart_learning_1986}. It involves calculating the gradients of the loss function relative to the network's weights and biases, which mathematically indicate the steepest increase in the loss function. Refer to Algorithm~\ref{Codes:GradientLossFunction} for a detailed overview of how to calculate these gradients. Then, gradient descent algorithm is used to iteratively update the weights and biases in the opposite direction of the gradients, thus reducing the loss and hopefully increasing the network's accuracy. Its variants, outlined with the respective formulas in Table~\ref{Tables:GradientDescent}, address different data and computational needs. Batch gradient descent uses the entire dataset for updates, ensuring completeness but potentially struggling with large datasets. Stochastic gradient descent updates the parameters per data point, offering speed at the cost of stability. Mini-batch gradient descent strikes a balance, using data subsets for more stable and computationally efficient updates \cite{goodfellow_deep_2016}.

\begin{algorithm}[htb!]
\caption{Gradient of the Loss Function \cite{goodfellow_deep_2016}.}
\label{Codes:GradientLossFunction}
\begin{algorithmic}
\REQUIRE \( f \): network function, \( \theta \): network parameters, \( x \): input data, \( y \): true label, \( l \): number of layers, \( L \): loss function
\STATE Compute the output of the network \( \hat{y} = f(x; \theta) \)
\STATE Compute the gradient of the loss function with respect to the output layer before activation:
\STATE \( \nabla_{z^{(l+1)}}L(f(x; \theta), y) = \frac{\partial L(f(x; \theta), y)}{\partial \hat{y}} \cdot \frac{\partial \hat{y}}{\partial z^{(l+1)}} \)
\FOR{\( \ell = l+1 \) to \( 1 \)}
\STATE Compute gradients of hidden layer parameters:
\STATE \( \nabla_{W^{(\ell)}}L(f(x; \theta), y) = \nabla_{z^{(\ell)}}L(f(x; \theta), y) \cdot h^{(\ell-1)T} \)
\STATE \( \nabla_{b^{(\ell)}}L(f(x; \theta), y) = \nabla_{z^{(\ell)}}L(f(x; \theta), y) \)
\STATE Compute gradient of previous hidden layer:
\STATE \( \nabla_{h^{(\ell-1)}}L(f(x; \theta), y) = W^{(\ell)T} \cdot \nabla_{z^{(\ell)}}L(f(x; \theta), y) \)
\STATE Compute gradient of previous hidden layer (before activation):
\STATE \( \nabla_{z^{(\ell-1)}}L(f(x; \theta), y) = \nabla_{h^{(\ell-1)}}L(f(x; \theta), y) \odot g'(z^{(\ell-1)}) \)
\ENDFOR
\end{algorithmic}
\end{algorithm}

In Table~\ref{Tables:GradientDescent}, the learning rate \( \eta \) is an hyperparameter that influences the convergence of gradient descent. An optimal learning rate balances convergence speed and precision, preventing overshooting or slow progress that could get trapped in local minima. Adaptive learning rate techniques by adjusting \( \eta \) during training improve learning efficacy, emphasising the balance between convergence speed and accuracy in neural network training. Some other techniques to improve convergence and its speed involve normalisation and standardisation of the input data \cite{goodfellow_deep_2016}.

\begin{table}[htb!]
\centering
\footnotesize
\begin{tabularx}{\textwidth}{@{}lXl@{}}
\toprule
\textbf{Algorithm} & \textbf{Formula} & \textbf{Description} \\
\midrule
Batch Gradient Descent & \(\theta := \theta - \eta \cdot \nabla_{\theta} J(\theta; X, y)\) & Updates parameters using all data \\
\addlinespace
Mini-Batch Gradient Descent & \(\theta := \theta - \eta \cdot \nabla_{\theta} J(\theta; X^{(i:i+n)}, y^{(i:i+n)})\) & Updates parameters using subsets of data \\
\addlinespace
Stochastic Gradient Descent & \(\theta := \theta - \eta \cdot \nabla_{\theta} J(\theta; x^{(i)}, y^{(i)})\) & Updates parameters per data point \\
\bottomrule
\end{tabularx}
\caption{Summary of the most common Gradient Descent variants \cite{goodfellow_deep_2016}.}
\label{Tables:GradientDescent}
\end{table}

\subsubsection{Bias-Variance Trade-off and Model Generalisation}

The generalisation power of a neural network depends on finding a balance between bias and variance. Bias leads to underfitting, where the model overlooks underlying trends due to overly simplistic assumptions, while variance leads to overfitting, where the model interprets noise as a signal due to excessive complexity. Both underfitting and overfitting result in poor performance of the model on new data, a significant issue in Forex trading, where accurate market signal detection is essential. To prevent this, cross-validation is commonly employed to evaluate a model's generalisation power in independent training, validation and testing datasets. Regularisation techniques such as L1 and L2 are also used to moderate the model complexity by adding penalties to the loss function. Furthermore, principal component analysis (PCA) can help to find the right balance between bias and variance by selecting, extracting and generating relevant features to the model. Implementing these strategies increases the change of training a model with more generalisation power, thus allowing it to perform in out-of-sample data \cite{goodfellow_deep_2016}. Although important, these techniques are not the main focus of this research and will not be explored in full detail.

\subsubsection{Universal Approximation Theorem}

The universal approximation theorem, formalised by Kurt Hornik in 1989, states \cite{hornik_multilayer_1989}:

\begin{quote}
"A Neural Network with one hidden layer and a linear output can approximate arbitrarily well any continuous function, given enough hidden units."
\end{quote}

This validates the capacity of neural networks to serve as universal function approximators, which is an important result in the context of creating predictive models for Forex trading because, in theory, they should be able to capture the complex underlying patterns that drive market movements, in conjunction with reinforcement learning, as we will see in the next sections.